\documentclass{beamer}
\mode<presentation>
\usepackage{amsmath}
\usepackage{amssymb}
%\usepackage{advdate}
\usepackage{graphicx}
\graphicspath{{./figs/}}
\usepackage{adjustbox}
\usepackage{subcaption}
\usepackage{enumitem}
\usepackage{multicol}
\usepackage{mathtools}
\usepackage{listings}
\usepackage{url}
\def\UrlBreaks{\do\/\do-}
\usetheme{Boadilla}
\usecolortheme{lily}
\setbeamertemplate{footline}
{
  \leavevmode%
  \hbox{%
  \begin{beamercolorbox}[wd=\paperwidth,ht=2.25ex,dp=1ex,right]{author in head/foot}%
    \insertframenumber{} / \inserttotalframenumber\hspace*{2ex} 
  \end{beamercolorbox}}%
  \vskip0pt%
}
\setbeamertemplate{navigation symbols}{}

\providecommand{\nCr}[2]{\,^{#1}C_{#2}} % nCr
\providecommand{\nPr}[2]{\,^{#1}P_{#2}} % nPr
\providecommand{\mbf}{\mathbf}
\providecommand{\pr}[1]{\ensuremath{\Pr\left(#1\right)}}
\providecommand{\qfunc}[1]{\ensuremath{Q\left(#1\right)}}
\providecommand{\sbrak}[1]{\ensuremath{{}\left[#1\right]}}
\providecommand{\lsbrak}[1]{\ensuremath{{}\left[#1\right.}}
\providecommand{\rsbrak}[1]{\ensuremath{{}\left.#1\right]}}
\providecommand{\brak}[1]{\ensuremath{\left(#1\right)}}
\providecommand{\lbrak}[1]{\ensuremath{\left(#1\right.}}
\providecommand{\rbrak}[1]{\ensuremath{\left.#1\right)}}
\providecommand{\cbrak}[1]{\ensuremath{\left\{#1\right\}}}
\providecommand{\lcbrak}[1]{\ensuremath{\left\{#1\right.}}
\providecommand{\rcbrak}[1]{\ensuremath{\left.#1\right\}}}
\theoremstyle{remark}
\newtheorem{rem}{Remark}
\newcommand{\sgn}{\mathop{\mathrm{sgn}}}
\providecommand{\abs}[1]{\left\vert#1\right\vert}
\providecommand{\res}[1]{\Res\displaylimits_{#1}} 
\providecommand{\norm}[1]{\lVert#1\rVert}
\providecommand{\mtx}[1]{\mathbf{#1}}
\providecommand{\mean}[1]{E\left[ #1 \right]}
\providecommand{\fourier}{\overset{\mathcal{F}}{ \rightleftharpoons}}
%\providecommand{\hilbert}{\overset{\mathcal{H}}{ \rightleftharpoons}}
\providecommand{\system}[1]{\overset{\mathcal{#1}}{ \longleftrightarrow}}
%\providecommand{\system}{\overset{\mathcal{H}}{ \longleftrightarrow}}
	%\newcommand{\solution}[2]{\textbf{Solution:}{#1}}
%\newcommand{\solution}{\noindent \textbf{Solution: }}
\providecommand{\dec}[2]{\ensuremath{\overset{#1}{\underset{#2}{\gtrless}}}}
\newcommand{\myvec}[1]{\ensuremath{\begin{pmatrix}#1\end{pmatrix}}}
\let\vec\mathbf

\lstset{
%language=C,
frame=single, 
breaklines=true,
columns=fullflexible
}

\numberwithin{equation}{section}

\title{2.10.11}
\author{ai25btech11022 - Narshitha}
% \maketitle
% \newpage
% \bigskip
\begin{document}
{\let\newpage\relax\maketitle}
\renewcommand{\thefigure}{\theenumi}
\renewcommand{\thetable}{\theenumi} 
\textbf{Question}:\\
A unit vector that is coplanar with $\hat{i}+\hat{j}+2\hat{k}$ and $\hat{i}+2\hat{j}+\hat{k}$ and perpendicular to $\hat{i}+\hat{j}+\hat{k}$ is\_\_\_\_\\
\textbf{Solution}:\\
Let the two vectors lie on plane $\vec{P}$\\
as origin also lies on the plane\\
The Equation of the plane is $\vec{r}.{\vec{n}^T}=0$
\begin{align}
    \myvec{1\\1\\2}.{\vec{n}}^T=0\\
    \myvec{1\\2\\1}.{\vec{n}}^T=0\\
    \myvec{1 && 1 && 2\\1 && 2 &&1}.\vec{n}=\myvec{0\\0}\\
\end{align}
The augmented  matrix is
\begin{align}
    \myvec{1 && 1 && 2 &&\vrule&& 0\\1 && 2 && 1 &&\vrule&& 0}\\
\end{align}  
$R_2\rightarrow R_2-R_1$ and $R_1\rightarrow R_1-R_2$\\
\begin{align}
    \myvec{1 && 0 && 3 && \vrule&& 0\\0 && 1 && -1 &&\vrule&& 0}
\end{align}
we get\\
\begin{align}
    \vec{n}=\myvec{-3\\1\\1}
\end{align}
Let the unit vector be $\vec{A}$\\

\begin{align}
\vec{B}=\myvec{1\\1\\1}\\
{\myvec{\vec{n} && \vec{B}}}^T.\vec{A}=\myvec{0\\0}\\
    \myvec{-3 && 1 && 1\\1 && 1 && 1}.\vec{A}=\myvec{0\\0}
\end{align}
The augmented matrix is
\begin{align}
    \myvec{-3 && 1 && 1 && \vrule &&0\\ 1 && 1 && 1 &&\vrule&& 0}
\end{align}
$R_1\longrightarrow R_1-R_2$ and $R_2 \longrightarrow R_2 + R_1 /4$
\begin{align}
    \myvec{-4 && 0 && 0 && \vrule &&0\\0&& 1 && 1 && \vrule&&0}\\
    \vec{A}=\myvec{0\\b\\-b}
\end{align}
As $\vec{A}$ is a unit vector
\begin{align}
 \therefore  \vec{A}=\myvec{0\\\frac{1}{\sqrt{2}}\\ -\frac{1}{\sqrt{2}}}
\end{align}
\frametitle{Graphical Representation}
\begin{center}
    \includegraphics[width=0.7\linewidth]{figs/fig1.png}
\end{center}
\end{document}
